% Structured abstract. Every number is from results/ at the commit named on the title page.
\begin{small}
\noindent\textbf{\sffamily Introduction.} Musculoskeletal deconditioning is one of the obstacles to long-duration
spaceflight, and head-down and horizontal bed rest are its established terrestrial analogues. Three decades of
bed-rest studies report lower-limb muscle loss, but they report it campaign by campaign, with different muscles,
durations, imaging methods and countermeasures. We set out to turn that literature into one documented dataset
and to build a modelling framework that estimates, and honestly validates, how much muscle is lost after a given
period of unloading and which muscles lose most.

\noindent\textbf{\sffamily Methods.} A systematic search of PubMed, Scopus, Web of Science and the NASA Technical
Reports Server (2013 onwards; 5,731 records), a pre-search corpus of older studies and one campaign from NASA's
open bed-rest archive were extracted into a frozen schema in which one row is one study, arm, muscle and
measurement occasion. Papers reporting the same participants were mapped to a single campaign, and the campaign,
not the paper, is the unit of every validation. The framework has three tiers: (1)~a three-level meta-regression,
nonlinear in the day of the scan, with random intercepts for campaign, study and muscle and cluster-robust
inference; (2)~four machine-learning families (ridge regression, random forest, support vector regression and
gradient boosting) compared with a duration-only curve under leave-one-campaign-out cross-validation with nested
tuning; and (3)~a post-hoc probabilistic forecast by a language model (TypeSafe's Jev) that reads each held-out
campaign's description alongside the other campaigns' data, checked by ablations, a recognition probe and a
validation battery. Every bar was declared before the answer it governs.

\noindent\textbf{\sffamily Results.} The frozen dataset holds 742 rows from 52 studies (1992--2026) and 36
independent campaigns; the modelling subset holds 346 rows from 32 campaigns, scanned between day~5 and day~119.
Loss follows a curved path: $-2.99$\pp{} per doubling of days (95\% CI $-3.44$ to $-2.53$), a time constant of 60
days and an eventual loss of $-17.3\%$ ($-20.9$ to $-13.7$) for the reference scenario, although logarithmic,
saturating and spline forms are indistinguishable (within 1.3 AIC points). Which muscle is measured matters most:
30\% of the variance lies between muscles within a campaign, plantar flexors lose 5.69\pp{} more than knee
extensors at day~60 ($p = 0.001$) and 5.60\pp{} more than dorsiflexors ($p = 0.003$), and countermeasure arms lose
3.73\pp{} less than controls ($p = 0.002$). No machine-learning family beat the duration curve (best: random forest,
3.18 vs.\ 3.16\pp{} out-of-campaign MAE), a null result pre-committed to in advance. The language-model forecast
was 0.42\pp{} more accurate than the curve (95\% CI 0.13 to 0.73; 20 of 32 campaigns), survived the removal of
every study-identifying detail and every presentation check, but missed its declared 15\% improvement bar and its
ranges were overconfident.

\noindent\textbf{\sffamily Conclusions.} The published bed-rest literature supports a quantified, curved
duration--response and a muscle-specific vulnerability led by the plantar flexors. At 32 independent campaigns,
flexible tabular learners add nothing to a simple curve; the sample size is the number of campaigns, not the
number of rows. The dataset and the framework are the contribution, and both are reproducible from one command.

\smallskip
\noindent\textbf{\sffamily Keywords:} bed rest; head-down tilt; muscle atrophy; spaceflight analogue;
meta-regression; leave-one-cohort-out cross-validation; machine learning; probabilistic forecasting.
\end{small}
